% TOP_PROBLEM: {"id": 30001752, "title": "Isomorphism of Free Group Factors", "category_id": 8, "category": {"id": 8, "name": "analysis", "display_name": "Analysis", "description": "Limits, continuity, calculus, and function theory.", "slug": "analysis", "order_index": 8, "created_at": "2026-07-31T15:26:25.671Z"}, "status": "open"}
% FIRST_POSTED: 2026-09-27
% !TeX program = lualatex
% Compile twice: lualatex 30001752.tex
% All references and prose are embedded; no BibTeX database or external inputs.
\documentclass[11pt,a4paper]{article}
\usepackage[margin=24mm,headheight=14pt]{geometry}
\usepackage{fontspec}
\setmainfont{Latin Modern Roman}
\setsansfont{Latin Modern Sans}
\setmonofont{Latin Modern Mono}
\usepackage{amsmath,amssymb,mathtools}
\usepackage{unicode-math}
\setmathfont{Latin Modern Math}
\usepackage{microtype}
\usepackage{enumitem,xurl,needspace}
\usepackage[dvipsnames]{xcolor}
\usepackage{hyperref}
\hypersetup{unicode=true,colorlinks=true,linkcolor=MidnightBlue,urlcolor=MidnightBlue,
 pdftitle={500 Mathematical Problem Statements},pdfauthor={Source-annotated compilation},bookmarksnumbered=true}
\usepackage{bookmark}
\usepackage{tocloft}
\setlength{\cftsecnumwidth}{3em}
\cftsetpnumwidth{2.5em}
\cftsetrmarg{4em}
\usepackage{titlesec}
\titleformat{\section}{\Large\bfseries\raggedright}{\thesection}{0.7em}{}
\usepackage{fancyhdr}
\pagestyle{fancy}\fancyhf{}
\fancyhead[L]{\small\sffamily Mathematical problem statements}
\fancyhead[R]{\small Source edition 22}
\fancyfoot[C]{\thepage}
\setlength{\parindent}{0pt}
\setlength{\parskip}{5pt plus 1pt}
\setlength{\emergencystretch}{3em}
\setcounter{tocdepth}{1}
\setcounter{secnumdepth}{1}
\setlist[itemize]{leftmargin=3em,itemsep=5pt,topsep=4pt,parsep=0pt}
\allowdisplaybreaks
\newcommand{\N}{\mathbb{N}}
\newcommand{\Z}{\mathbb{Z}}
\newcommand{\Q}{\mathbb{Q}}
\newcommand{\R}{\mathbb{R}}
\newcommand{\C}{\mathbb{C}}
\newcommand{\F}{\mathbb{F}}
\newcommand{\A}{\mathbb{A}}
\newcommand{\PP}{\mathbb P}
\newcommand{\E}{\mathbb{E}}
\newcommand{\eps}{\varepsilon}
\newcommand{\dd}{\,\mathrm d}
\newcommand{\sref}[2]{\hyperlink{source:#1:#2}{[S#2]}}
\newcommand{\eref}[2]{\hyperlink{extra:#1:#2}{[E#2]}}
% Turn the next switch off for a shorter reading copy. The quotations preserve
% every exactTarget field, including qualifications omitted from a short summary.
\newif\ifshowcatalogue
\showcataloguetrue
\newcommand{\cataloguescope}[1]{\ifshowcatalogue
 \paragraph{Exact scope in the supplied catalogue.}
 \begin{quote}\small #1\end{quote}\fi}

% Standalone notebook wrapper; original problem section retained verbatim.
\numberwithin{equation}{section}
\hypersetup{pdftitle={Research attempt -- problem 30001752},
 pdfauthor={Research notebook; source compilation retained}}
\fancyhead[L]{\small\sffamily Research notebook}
\fancyhead[R]{\small\sffamily Problem 30001752 | 27 September 2026}
\begin{document}
\setcounter{section}{0}
% ================================================================
% problemId: 30001752
% Canonical problem ID: 30001752
% exactTarget SHA-256: 799cc08523718dbc763a4beff512e44304f860fa456e8beba4962dd069aeb9fe
\clearpage
\section*{\texorpdfstring{Isomorphism of Free Group Factors}{Isomorphism of Free Group Factors}}\label{30001752}
\subsection{Definitions and mathematical statement}
For a discrete group $G$, the left regular representation on $\ell^2(G)$ is $\lambda(g)\delta_h=\delta_{gh}$. Its group von Neumann algebra is
$L(G)=\{\lambda(g):g\in G\}''$, the double commutant in bounded operators. Let $F_m$ be the free group on $m$ generators. Determine the isomorphism relation among $L(F_m)$ for distinct ranks in the nonabelian range, particularly
\[
 L(F_2)\stackrel{?}{\cong}L(F_3).
\]
Isomorphism here means a von Neumann algebra $*$-isomorphism, not an isomorphism between the underlying free groups. For these finite factors the normalized trace is intrinsic. The issue is whether the operator algebra remembers the number of generators.
\par\smallskip\textit{Formulation and concept references:} \sref{30001752}{1}.
\cataloguescope{Determine whether L(F\_m) and L(F\_n) are isomorphic when m != n, especially whether L(F\_2) is isomorphic to L(F\_3).}
\subsection{Short English statement}
Can free groups with different numbers of generators produce isomorphic von Neumann algebras, especially at ranks two and three?
% BEGIN RESEARCH ADDITION ATTEMPT-174
\subsection{Research attempt}

\paragraph{Current frontier.}
The random-matrix strong-convergence results discussed in \sref{30001752}{1},
Section 2.2, approximate prescribed free generators; they do not identify
all possible generating families of their von Neumann algebra. Free entropy
dimension distinguishes the standard free semicircular families, but the
needed invariance under arbitrary von Neumann algebra generation is the
obstacle to using it for the stated isomorphism problem
\eref{30001752}{2}, Sections 2.4--2.5. Jung's packing formulation makes that
obstacle quantitative \eref{30001752}{1}, Definition 2.1 and Corollary 2.4.
The argument below investigates this route without identifying group rank
with an already established factor invariant.

\paragraph{Attack route.}
A compression argument would require an isomorphism invariant with a known
nontrivial scaling law. A generator-count argument must be invariant under
changing to an arbitrary generating tuple, not just a polynomial change of
coordinates. We pursue the metric-entropy version of the second route:
quantify how accurately and how gently one must reconstruct one microstate
tuple from another. The missing lemma will be a uniform, subpower-distortion
transport statement. A scalar polynomial example tests the proposed shortcut
from approximation alone to that statement.

\paragraph{Attempt: a finite-matrix packing lemma.}
For tuples of self-adjoint $k\times k$ matrices use
\[
 d_k(A,B)^2=\sum_i k^{-1}\operatorname{Tr}((A_i-B_i)^2).
\]
Let $\mathcal X_k,\mathcal Y_k$ be bounded sets of such tuples, possibly
of different lengths, and let $\Phi_k:\mathcal X_k\to\mathcal Y_k$.
Suppose there is a map $\Psi_k$ defined on $\Phi_k(\mathcal X_k)$ with
\[
 d_k(\Psi_k\Phi_k(A),A)\le\eta,
 \qquad d_k(\Psi_k(B),\Psi_k(B'))\le Ld_k(B,B').
\]
Let $\operatorname{sep}_\eps$ denote the largest cardinality of a set whose
pairwise distances exceed $\eps$. If $2\eta<\eps$, then
\begin{equation}\label{eq174-packing}
 \operatorname{sep}_\eps(\mathcal X_k)
 \le\operatorname{sep}_{(\eps-2\eta)/L}(\mathcal Y_k).
\end{equation}
Indeed, the triangle inequality bounds
$d_k(A,A')$ by $2\eta+L d_k(\Phi_k(A),\Phi_k(A'))$.
Thus images of an $\eps$-separated set remain separated at the claimed
scale. Injectivity on that set follows at once. This proof has no dependence
on $k$ and requires no differentiability or volume-preservation assumption.

On a fixed operator-norm ball of radius $R\ge1$, a noncommutative
polynomial $P=\sum_w c_w w$ is Lipschitz in this metric, with a bound
obtained by summing $|c_w||w|R^{|w|-1}$, up to the fixed tuple-norm
constant. To check this, telescope the difference of each monomial, one
factor at a time, and use $\|ABC\|_2\le\|A\|\|B\|_2\|C\|$.
Thus polynomial reconstruction is suitable for \eqref{eq174-packing},
but its coefficient sizes and degrees matter.

\paragraph{Attempt: the precise transport hypothesis.}
For a self-adjoint tuple $X$, let $\Gamma_R(X;m,k,\gamma)$ consist of
operator-norm-$R$ matrix tuples whose normalized traces of words of length
at most $m$ differ from those of $X$ by less than $\gamma$.
Define the normalized packing quantity by taking $\limsup_{k\to\infty}$
of $k^{-2}\log\operatorname{sep}_\eps$, then the infimum over
$(m,\gamma)$. Denote it by $\mathcal P_{\eps,R}(X)$.
The packing description of $\delta_0$ takes a further small-scale limit;
fixed changes in the separation convention do not change the dimension
\eref{30001752}{1}, Corollary 2.4.

Here is a sufficient input for transporting that dimension from $X$ to
$Y$. Fix cutoffs $R_X,R_Y$ larger than the operator norms of the tuples.
For every small $\eps$ suppose that, for \emph{every} target moment
condition $(m_Y,\gamma_Y)$, sufficiently stringent source moment conditions
$(m_X,\gamma_X)$ admit maps as in \eqref{eq174-packing} between the
corresponding microstate sets, for all sufficiently large $k$. Require
\begin{equation}\label{eq174-distortion}
 \eta\le\eps/4,\qquad L=L(\eps)\ge1,
 \qquad \log L(\eps)=o(\log(1/\eps)).
\end{equation}
The constants and cutoffs are independent of the matrix dimension and the
target moment accuracy. The maps themselves may depend on those moment
conditions. Then \eqref{eq174-packing}, followed by the indicated limits
and infima, gives
\begin{equation}\label{eq174-entropy}
 \mathcal P_{\eps,R_X}(X)
       \le \mathcal P_{\eps/(2L(\eps)),R_Y}(Y).
\end{equation}
The order of quantifiers is essential: a map into one fixed coarse target
microstate set would not justify taking the target infimum.

Divide by $\log(1/\eps)$. By \eqref{eq174-distortion},
\[
 \frac{\log(2L(\eps)/\eps)}{\log(1/\eps)}\longrightarrow1.
\]
The usual fixed-cutoff packing formula therefore yields
$\delta_0(X)\le\delta_0(Y)$. A reverse transport with the same hypotheses
yields equality. This proves the conditional dimension-transfer statement;
there is no exchange of the infimum over moment accuracy with a matrix
limit that was not part of the definition.

In particular, applying this transport to every pair of generating tuples
of an isomorphic pair of free group factors would make $\delta_0$ a factor
invariant in the needed cases. Its value on a free semicircular $r$-tuple
is $r$, so this would distinguish ranks two and three, and more generally
distinct finite nonabelian ranks; see \eref{30001752}{2} for the free-entropy
calculation and its intended application. This is a sufficient deep input,
not a property established merely from an isomorphism of factors.

\paragraph{Attempt: breaking the approximation-to-transport shortcut.}
Approximation of actual operators does not automatically control its
Lipschitz cost. Even on $[-1,1]$, put
\[
 P_n(x)=x+n^{-1}T_{n^2}(x),
\]
where $T_j(\cos\theta)=\cos(j\theta)$. Then
$\|P_n-\mathrm{id}\|_\infty\le1/n$, while
\[
 P_n'(1)=1+n^{-1}(n^2)^2=1+n^3.
\]
The derivative identity follows by taking the limit of
$j\sin(j\theta)/\sin\theta$ as $\theta\to0$.
For an approximation scale $\eps=1/n$, this Lipschitz cost is a positive
power of $1/\eps$, not the subpower cost in \eqref{eq174-distortion}.
If $L\asymp\eps^{-c}$ in \eqref{eq174-entropy}, the elementary dimension
comparison can lose the factor $1+c$ instead of proving equality.

This example does not say that every approximation to the identity must be
bad: the identity itself is an excellent choice. It disproves the general
inference that \emph{an} accurate sequence of polynomial approximants comes
with adequate quantitative control. Von Neumann generation supplies strong
or $L^2$ approximation, not the fixed-cutoff, uniformly reversible,
subpower-distortion microstate transport asserted above. Controlling only
one actual operator tuple also does not control every matrix tuple with
sufficiently many matching moments unless the relevant estimates are
transferred with their tolerances.

\paragraph{Stress test.}
The normalized Hilbert--Schmidt metric, rather than unnormalized Euclidean
matrix distance, is used throughout. A dimension-dependent Lipschitz
constant could survive division by $k^2$ and invalidate the argument.
Operator-norm cutoffs cannot be allowed to diverge arbitrarily with the
small scale without checking the entropy definition again.

The entropy route studies microstates with moment tolerances. Strong
convergence of a selected random-matrix sequence \sref{30001752}{1} is valuable
but is not a uniformly invertible parametrization of those entire spaces.
Nor would a group-theoretic invariant of $F_r$ automatically be an invariant
of $L(F_r)$: isomorphisms need not carry canonical group unitaries to group
unitaries. The argument makes no assertion about countably infinite rank
from the finite-tuple calculation alone.

\paragraph{Outcome.}
\textbf{Outcome: Conditional reduction.}

The attempt proves a dimension-free packing inequality and a quantified
sufficient condition for invariance of free entropy dimension under change
of generators. A concrete polynomial test shows why approximation alone
cannot supply the condition. The unproved input is the uniform transport
statement for arbitrary generating tuples; obtaining it in the free-factor
cases would resolve the finite-rank distinction. No factor isomorphism or
nonisomorphism is established here, and no novelty is claimed for the
standard packing interpretation itself.
% END RESEARCH ADDITION ATTEMPT-174
\subsection{Sources}
\begin{itemize}\raggedright
\item[\textnormal{[S1]}]\hypertarget{source:30001752:1}{} Ramon van Handel, Strong Convergence: A Short Survey (2025), section 2.2.
\url{https://arxiv.org/html/2510.12520v1}.
{\small\textit{Catalogue formulation source.}}
% BEGIN RESEARCH ADDITION SOURCES-174
\item[\textnormal{[E1]}]\hypertarget{extra:30001752:1}{} Kenley Jung,
\emph{A free entropy dimension lemma}, Pacific Journal of Mathematics
211(2) (2003), 265--271; arXiv:math/0207149v2,
Definition 2.1, Remark 2.3, and Corollary 2.4.
\url{https://arxiv.org/html/math/0207149}.
{\small\textit{Consulted 27 September 2026; normalized microstate packing
and the order of its limits.}}
\item[\textnormal{[E2]}]\hypertarget{extra:30001752:2}{} Dan Voiculescu,
\emph{Free Entropy}, arXiv:math/0103168, 2001,
Sections 2.4--2.5 and the generator-invariance problem.
\url{https://arxiv.org/html/math/0103168}.
{\small\textit{Consulted 27 September 2026; free-family dimension and the
distinction between a tuple invariant and a von Neumann algebra invariant.}}
% END RESEARCH ADDITION SOURCES-174
\end{itemize}

\end{document}
